\begin{helper}
Let \(N,n,m\) be natural numbers with \(m\le N\).  Put
\[
q=\frac{m}{N}
\]
using Lean's real-division convention when \(N=0\), and put \(\mu=nq\).  For
every \(L\in\mathbb R\),
\[
(1-q+q e^L)^n\le \exp(\mu(e^L-1)).
\]
\end{helper}
\begin{proof}
The real number \(q\) is nonnegative because it is the quotient of two
nonnegative natural-number casts.  Also \(q\le 1\).  If \(N=0\), then
\(m\le N\) forces \(m=0\), so \(q=0\).  If \(N>0\), then division by the
positive real number \(N\) preserves the inequality \(m\le N\), giving
\(m/N\le 1\).

Since \(0\le q\le 1\) and \(e^L>0\), both terms \(1-q\) and \(q e^L\) are
nonnegative, hence \(1-q+q e^L\ge0\).  The elementary exponential inequality
\(1+x\le e^x\), applied to \(x=q(e^L-1)\), gives
\[
1-q+q e^L = 1+q(e^L-1)\le \exp(q(e^L-1)).
\]
Because the left base is nonnegative, raising this inequality to the natural
power \(n\) preserves it:
\[
(1-q+q e^L)^n\le \exp(q(e^L-1))^n.
\]
Finally \(\exp(a)^n=\exp(na)\) for natural \(n\), and with
\(a=q(e^L-1)\) this right hand side is
\[
\exp(nq(e^L-1))=\exp(\mu(e^L-1)).
\]
\end{proof}
